% Public source snapshot of the instructor's Module 11 lecture notes. Executable
% implementations, diagram source, external images, private links, and product-current
% claims are omitted under the boundary in sources/README.md. Headings and equations
% are preserved for provenance, not endorsed as reproduced results; the book
% independently re-derives the mathematics and checks claims against primary sources.
\documentclass{article}
\usepackage{amsmath}
\usepackage{amssymb}
\usepackage[margin=1in]{geometry}
\usepackage{xcolor}
\usepackage{tcolorbox}
\usepackage{booktabs}
\usepackage{hyperref}

\title{Lecture 11.4: Reinforcement Learning for Fine-Tuning}
\author{Deep Learning Class}
\date{\today}

\begin{document}
\maketitle

\begin{tcolorbox}[colback=gray!5!white, colframe=gray!60!black, title=Public Snapshot Boundary]
This provenance snapshot preserves the instructor-authored lecture spine and equations.
Executable and external-library implementations, diagram source, external images,
private links, and claims tied to current products have been omitted; their original
locations remain marked in the source. Remaining technical and numerical claims are
historical lecture notes and are independently checked before publication in the book.
\end{tcolorbox}

\section{Reinforcement Learning for Fine-Tuning (Brief Overview)}

\subsection{Motivation: Why Not Just Supervised Fine-Tuning?}

Supervised Fine-Tuning (SFT) trains models to mimic demonstrations. Given input-output pairs $(x, y)$, we minimize:

\begin{equation}
\mathcal{L}_{\text{SFT}} = -\mathbb{E}_{(x,y) \sim \mathcal{D}} \left[ \log p_\theta(y \mid x) \right]
\end{equation}

\textbf{Problem:} SFT teaches the model ``what good outputs look like'' but not ``what makes an output better than another.''

\begin{tcolorbox}[colback=yellow!5!white, colframe=orange!50!black, title=SFT Limitations]
\textbf{Example:} You want a helpful, harmless assistant.

\textbf{SFT approach:} Train on demonstrations of helpful responses.

\textbf{Issue:} Model learns the \textit{format} of helpfulness but may still generate:
\begin{itemize}
\item Plausible but incorrect answers (hallucinations)
\item Responses that are technically correct but unhelpful
\item Content that follows the style but violates safety guidelines
\end{itemize}

\textbf{What's missing?} Explicit optimization for desirable properties (helpfulness, truthfulness, safety).
\end{tcolorbox}

\subsection{High-Level Idea: Reinforcement Learning from Human Feedback (RLHF)}

Instead of only imitating demonstrations, we can:

\begin{enumerate}
\item \textbf{Learn what humans prefer} (train a reward model)
\item \textbf{Optimize the language model} to generate high-reward outputs
\end{enumerate}

\textbf{Three-stage pipeline:}

\begin{center}
% [Diagram source omitted from the public snapshot; the book uses independently drawn figures.]
\end{center}

\textbf{Stage 1 (SFT):} Train base model on demonstrations $\rightarrow$ produces $\pi_{\text{SFT}}$

\textbf{Stage 2 (Reward Modeling):} Collect human preferences on pairs of outputs $(y_1, y_2)$ for the same prompt $x$. Train a reward model $r_\phi(x, y)$ to predict which output humans prefer:

\begin{equation}
\mathcal{L}_{\text{RM}} = -\mathbb{E} \left[ \log \sigma\left(r_\phi(x, y_{\text{win}}) - r_\phi(x, y_{\text{lose}})\right) \right]
\end{equation}

where $\sigma$ is the sigmoid function and $y_{\text{win}}$ is the preferred output.

\textbf{Stage 3 (RL Optimization):} Optimize the language model policy $\pi_\theta$ to maximize expected reward while staying close to the SFT model (to prevent collapse):

\begin{equation}
\mathcal{L}_{\text{RL}} = \mathbb{E}_{x \sim \mathcal{D}, y \sim \pi_\theta} \left[ r_\phi(x, y) - \beta \cdot D_{\text{KL}}\left(\pi_\theta(y|x) \| \pi_{\text{SFT}}(y|x)\right) \right]
\end{equation}

The KL penalty term prevents the model from drifting too far from the SFT initialization.

\begin{tcolorbox}[colback=blue!5!white, colframe=blue!50!black, title=Intuition Without RL Theory]
Think of the reward model as a ``critic'' that scores outputs. The RL stage adjusts the language model to generate outputs that the critic rates highly, but with a constraint: don't change so much that you lose the knowledge from SFT.

\textbf{You don't need to understand PPO internals} to follow this conceptual overview.
\end{tcolorbox}

\subsection{Direct Preference Optimization (DPO): Simpler Alternative}

\textbf{Problem with RLHF:} Requires training a separate reward model and then running RL (complex, unstable).

\textbf{DPO insight:} You can optimize directly on preference pairs without an explicit reward model!

Given preference data $\{x, y_{\text{win}}, y_{\text{lose}}\}$, DPO minimizes:

\begin{equation}
\mathcal{L}_{\text{DPO}} = -\mathbb{E} \left[ \log \sigma\left( \beta \log \frac{\pi_\theta(y_{\text{win}}|x)}{\pi_{\text{ref}}(y_{\text{win}}|x)} - \beta \log \frac{\pi_\theta(y_{\text{lose}}|x)}{\pi_{\text{ref}}(y_{\text{lose}}|x)} \right) \right]
\end{equation}

where $\pi_{\text{ref}}$ is the reference policy (usually the SFT model).

\textbf{Advantage:} No reward model, no RL algorithm. Just a classification-like loss on preference pairs.

\textbf{When to use DPO:}
\begin{itemize}
\item You have preference data (A vs B comparisons)
\item You want simplicity over the full RLHF pipeline
\item You're working with smaller models or limited compute
\end{itemize}

\subsection{External Library Workflow}

The original lecture included a survey of one external library's trainer APIs.
That product-current API inventory is omitted from this public snapshot.

\textbf{Typical external-library workflow:}

% [Executable library example omitted from the public snapshot: independent provenance was not established. See sources/README.md.]

The original lecture demonstrated this workflow with an external library. That executable example is omitted from this provenance snapshot.

\subsection{When to Use Each Method}

\begin{center}
\begin{tabular}{|l|p{10cm}|}
\hline
\textbf{Method} & \textbf{Use When...} \\
\hline
SFT & You have demonstrations; task is well-defined; format matters \\
\hline
RLHF (PPO) & You have human preferences; want to optimize complex objectives (helpfulness + safety); have compute for reward model + RL \\
\hline
DPO & You have preference pairs; want simpler pipeline; smaller models \\
\hline
\end{tabular}
\end{center}

\begin{tcolorbox}[colback=green!5!white, colframe=green!50!black, title=Practical Recommendation]
\textbf{For this course:} Start with SFT (Sections 3-5). If you need alignment:
\begin{enumerate}
\item Collect preference data (have humans rank outputs)
\item Try DPO first (simpler)
\item Use RLHF only if DPO is insufficient
\end{enumerate}

% [Product-current adoption claim omitted; publication claims are checked against primary sources.]
\end{tcolorbox}

\subsection{Example: DPO Workflow (Code Omitted)}

The original lecture pointed to a separate implementation file. That reference and the uncleared executable example are omitted from this public snapshot.

% [Executable library example omitted from the public snapshot: independent provenance was not established. See sources/README.md.]

\textbf{Key points:}
\begin{itemize}
\item The reference policy is frozen (usually the SFT model)
\item The active policy is optimized to prefer ``chosen'' over ``rejected''
\item $\beta$ controls how much the policy can deviate from the reference
\end{itemize}

\subsection{What We're Not Covering (But You Can Learn More)}

This course focuses on deep learning fundamentals, not RL theory. We're \textit{not} covering:

\begin{itemize}
\item Markov Decision Processes (MDPs)
\item Policy gradient derivations
\item Value functions and Q-learning
\item PPO algorithm internals (clipped objective, advantage estimation)
\item Variance reduction techniques
\end{itemize}

\textbf{If you want to learn RL for LLMs:}
\begin{itemize}
\item \textbf{TRL Documentation:} \url{https://huggingface.co/docs/trl}
\item \textbf{RLHF Paper:} Ouyang et al., ``Training language models to follow instructions with human feedback'' (2022)
\item \textbf{DPO Paper:} Rafailov et al., ``Direct Preference Optimization'' (2023)
\item \textbf{Spinning Up in Deep RL:} \url{https://spinningup.openai.com} (for RL foundations)
\end{itemize}

\subsection{Summary: The Alignment Stack}

\begin{center}
% [Diagram source omitted from the public snapshot; the book uses independently drawn figures.]
\end{center}

\textbf{For this course:} Focus on SFT (Sections 3-5). Treat RLHF/DPO as ``optional advanced topics.''

% [Product-current library recommendation omitted from the public snapshot.]

\end{document}
